At any given interest rate, an increase in capital of 5000 Thl. can cover a pension increase of 10 Thl. for only a certain number of pensioners; at an interest rate of 3 percent\footnote{When the proposal for the regulation under discussion was submitted in July 1794, the calculation was based, in unambiguous words, only on this interest rate and not on 4 percent. Accordingly, the following passage in the paper by University Councillor O. dated 15 December 1844, p. 5:
\begin{quote}\raggedright
‘The provision made in 1794 .\,.\,.\,.\,. was probably based on a calculation\par
‘that the interest on the fund raised to 5000 Thl., at 4 percent, would amount to 200 Thl.,\par
‘of which \(\tfrac14\) was to be added back to the capital, while \(\tfrac34\) was to be divided among the widows,\par
‘which, with 15 widows, would give 10 Thl. to each’
\end{quote}
must be corrected. There is \textit{nothing whatsoever} about such a calculation, about interest at 4 percent, or about setting aside a quarter of it in the deliberations of 1794, which will shortly be examined more closely.} for at most
\SourceLeaf{136}
15, at an interest rate of \(3\tfrac12\) percent for at most 17, and at an interest rate of 4 percent for at most 20 pensioners; when these numbers are reached, the interest yield of that increase in capital is entirely absorbed by them. Thus, if the number of those to receive pensions rises above these respective limits, there will be nothing to meet all the remaining pensions except 1) the surplus which the fund's annual income had provided before the increases in capital and pension rates, or rather would have provided if the respective standard number of pensions (15; 17; 20) had existed at that time\footnote{In fact, the income at that time (1794) had not quite been sufficient even for 18 pensions, although an annual subsidy of 150 Thl. was then still being paid from the church treasury, which later ceased. The above estimate is, however, subject to a modification, since, chiefly as a result of irregularities during the Westphalian government, the rule of progression was not followed exactly. Had it been followed with complete strictness, the pension rates since 1815 would always have had to be raised when the capital was still lower than it actually was when they were raised, — and the present capital assets would perhaps be 5000 Thl. poorer.}. 2) Any increase in the pharmacy rent brought about since that time. 3) The increased income from the members' contributions owing to their greater number. But these uncertain and meagre resources would soon be exhausted if those standard numbers were exceeded further, and all the sooner the higher the pension rate itself had already risen. This last remark is of great importance inasmuch as it shows with the utmost clarity that the danger arising from consistent adherence to the second interpretation becomes all the greater the more pension increases have already taken place before the standard number is exceeded.

Although it was remarked above that no proviso can be found in the printed regulations by which the fund could be prevented from bleeding to death sooner or later from this source, I must not conceal that the receipt forms on which the widows collect their pensions include the proviso that the pension is fixed at .\,.\,.\,.\,. Rthl. for the present, and \textit{for as long as the fund's circumstances will permit}. Probably few members of the widows' fund know this clause: at least I myself, although I have been a participant for 38 years, became aware of it only very recently. It is not, however, intended, as I initially supposed, to guard against possible dangers to be feared from the rule of progression; for it already appears, in precisely the same words, in the printed receipt forms before 1794. Such provisos, framed in entirely general terms, may have their merits in some respects: in the ordinary course of events, however, unless extraordinary circumstances arise, they remain unused until the most urgent necessity compels their application. No criterion for a decision, no standard is given by which one can recognise whether the fund's circumstances permit payment or do not permit it. If one waits until it has already repeatedly become necessary to draw on capital as well to pay the pensions, one has very probably already waited much too long, and what was remarked at [p. 128, l. 12] will then prove true. In any case, with regard to maintaining the fund's reputation for soundness, there is a great difference between a fall in the pension rate resulting from a quite definite, fixed rule (as on the first interpretation), and a reduction of pensions which the fund finally finds itself compelled to make because it can no longer meet the obligations which, under the
\SourceLeaf{137}
statute (on the second interpretation), it has in fact undertaken \textit{unconditionally}.

